% ============================================================
% NEW §VI: Discussion
% Insert before Conclusion (currently §VI → becomes §VII)
% ============================================================

\section{Discussion}
\label{sec:discussion}

\subsection{Generality of the Framework}

The time-aware evaluation framework presented in this paper is deliberately
detector-agnostic and task-agnostic.  Three design choices enable this generality.

\textbf{Decoupling from the sliding-window assumption.}
DR@$\Delta t$, ADD, and MTBFA are defined over sequences of timestamped binary
decisions $\hat{y}_k$ and require only that (i) each decision carries a causal
timestamp and (ii) an event onset is defined.
Detectors that issue per-step decisions (e.g.\ online RNNs, event-triggered
architectures) can be evaluated by the same formulae without modification.

\textbf{Portability to other event-detection domains.}
The three metrics apply unchanged to any system that must raise a first alarm
within a budget, where repeated false positives are operationally disruptive.
Examples include industrial fault prediction (alarm management standards require
MTBFA $\geq 10\,\text{min}$ \cite{eemua191}), medical monitoring (response
within the golden hour \cite{iso11064}), and network intrusion detection where
analyst response-time constraints are well-documented \cite{axelsson2000base}.
Adapting the framework requires only re-grounding the numerical thresholds to the
target domain's physical risk model.

\textbf{Framework vs.\ case study.}
The seven detectors evaluated in Section~\ref{sec:casestudy} serve as a
\emph{case study} demonstrating the framework's discriminative power, not as
an endorsement of any particular architecture.
The specific metric values will change with training data, hyperparameters,
and the spoofing environment; what remains invariant is the framework's ability
to reveal detectors that pass conventional metrics but fail operational
time-aware standards.

\subsection{Limitations}

\textbf{Simulation-to-reality gap.}
Our primary dataset is generated by the MATLAB UAV Toolbox, whose GPS sensor
model uses additive Gaussian white noise.
Real GNSS receivers exhibit correlated noise from multipath propagation,
ionospheric delay, and receiver clock drift—phenomena not captured by the simulator.
The zero-shot cross-domain evaluation (Section~\ref{sec:cross_domain}) quantifies
the resulting performance degradation and indicates the gap is moderate,
but large-scale in-the-wild deployment validation remains future work.

\textbf{Threshold sensitivity to mission profile.}
The DR@5s and MTBFA thresholds are derived for power-line inspection at 5 m/s.
For missions with different speeds, safe-distance requirements, or operator
workload budgets, the thresholds require re-derivation;
Table~\ref{tab:deployability} therefore reports the full DR@$\{1,3,5,10,15\}\,\text{s}$
and MTBFA values so that practitioners can apply their own thresholds.

\textbf{Bias toward abrupt attacks in the sliding-window setting.}
End-point labelling assigns the window label from the final sample.
For \emph{slow drift} attacks, the earliest windows may straddle the onset,
delaying the first positive label and potentially inflating ADD.
This limitation is noted in Section~\ref{sec:window_labeling} and is
consistent with how operator-observable disturbances manifest in practice
(slow spoofing is imperceptible to the navigation filter for the first few seconds).

\textbf{Single simulation platform.}
All primary experiments use one simulation toolchain.
Future work should cross-validate with FlightGear, Gazebo, or X-Plane to assess
inter-simulator variance, and with additional real-world hardware platforms
beyond the datasets used here.
